%% ========================================================================
%% Lampiran A: Referensi REST API
%% ========================================================================
%% STATUS: SCAFFOLD ONLY - lihat docs/authoring-guide.md
%% ========================================================================

\chapter{Referensi REST API}
\label{app:rest-api}

Lampiran ini merangkum endpoint REST API Simple POS yang dipakai oleh klien
kasir mobile, sebagaimana dibahas pada \cref{ch:membangun-rest-api}. Seluruh
endpoint diawali dengan Sanctum sebagai lapisan autentikasi berbasis token.

\section{Autentikasi}

\begin{tblr}{
  colspec = {X[0.4,l] X[1.4,l] X[1.4,l]},
  hline{1,2,Z} = {solid},
}
\textbf{Method} & \textbf{Endpoint} & \textbf{Keterangan} \\
POST & \route{/api/login} & Mengirim email dan password, mengembalikan token Sanctum \\
POST & \route{/api/logout} & Mencabut token yang sedang dipakai (memerlukan header \route{Authorization}) \\
\end{tblr}

\begin{notebox}
Setiap permintaan ke endpoint yang dilindungi wajib menyertakan header
\route{Authorization: Bearer <token>}, dengan \route{<token>} adalah nilai
yang dikembalikan oleh \route{POST /api/login}.
\end{notebox}

\section{Produk}

\begin{tblr}{
  colspec = {X[0.4,l] X[1.4,l] X[1.4,l]},
  hline{1,2,Z} = {solid},
}
\textbf{Method} & \textbf{Endpoint} & \textbf{Keterangan} \\
GET & \route{/api/products} & Mengambil daftar produk beserta stok dan harga \\
\end{tblr}

\section{Transaksi}

\begin{tblr}{
  colspec = {X[0.4,l] X[1.4,l] X[1.4,l]},
  hline{1,2,Z} = {solid},
}
\textbf{Method} & \textbf{Endpoint} & \textbf{Keterangan} \\
POST & \route{/api/transactions} & Membuat transaksi baru; total dihitung server-side \\
GET & \route{/api/transactions} & Mengambil daftar transaksi milik pengguna yang sedang masuk \\
GET & \route{/api/transactions/\{id\}} & Mengambil detail satu transaksi \\
\end{tblr}

\begin{warningbox}
Sama seperti pada versi web, \route{POST /api/transactions} tidak pernah
menerima nilai total dari klien. Total tetap dihitung ulang di server dari
harga produk dan kuantitas, mengikuti prinsip yang sama seperti pada
\cref{ch:validasi-input}.
\end{warningbox}
